\chapter{Objetivo do manual}
Este manual foi escrito para RH, chefias com permissão de gestão e perfis equivalentes. O objetivo é documentar a operação real da plataforma com detalhe suficiente para que a pessoa saiba não apenas \textit{o que existe}, mas \textit{como usar} e \textit{quando usar} cada funcionalidade.

\chapter{Visão geral da plataforma}
A Smarter Hub reúne os principais processos de pessoas: colaboradores, permissões, férias, ausências, formações, admissões, banco de horas, notificações, saúde e bem-estar e dashboards de acompanhamento.

\section{Princípio de funcionamento}
Cada módulo tenta manter o fluxo dentro do portal. Em vez de usar vários canais paralelos, RH trabalha a partir de listas, detalhe, modal, aprovação, exportação e histórico de ações.

\section{Nomenclatura de trabalho}
\begin{longtable}{|p{0.34\textwidth}|p{0.58\textwidth}|}
\hline
\textbf{PT-PT} & \textbf{PT-BR / observação} \\
\hline
Chefes de equipa & Líderes de equipe / gestores de equipe. \\
\hline
A Minha Ficha & Meu cadastro / meu perfil, dependendo do contexto da conversa. \\
\hline
Formações & Treinamentos. \\
\hline
Férias / Ausências & Mantém-se o nome, mas os motivos e a leitura administrativa podem variar por país. \\
\hline
Saúde e bem-estar & Mantém-se o nome, com conteúdos comuns e conteúdos locais. \\
\hline
\end{longtable}

\chapter{Como trabalhar na plataforma}
\section{Entrada e contexto}
Ao entrar, verifique logo o seu contexto: país, equipa, permissões e área visível. Em RH é particularmente importante perceber se está em modo de consulta, edição, aprovação ou exportação.

\section{Boas práticas operacionais}
\begin{itemize}
  \item Confirme sempre o país do colaborador antes de editar ou aprovar.
  \item Leia a página até ao fim antes de gravar.
  \item Utilize filtros de pesquisa para evitar abrir registos errados.
  \item Feche modais ou detalhes depois de validar a informação.
  \item Use exportações como apoio e não como substituto da validação de detalhe.
\end{itemize}

\chapter{Dashboard}
\section{Dashboard}
O Dashboard é a vista executiva da organização. Está orientado para análise, acompanhamento e exportação de informação.

\subsection{O que normalmente existe}
\begin{itemize}
  \item Filtros de hierarquia, geografia, género, função ou outros eixos analíticos.
  \item Cartões e métricas agregadas.
  \item Tabelas de detalhe.
  \item Drill-down e drill-through.
  \item Exportação para Excel.
  \item Modalizações de detalhe, como o caso do voucher NOS quando aplicável.
\end{itemize}

\subsection{Como usar corretamente}
\begin{enumerate}
  \item Abra o Dashboard.
  \item Escolha a dimensão que pretende analisar.
  \item Ajuste os filtros até isolar o universo desejado.
  \item Use o drill-up ou drill-down para navegar entre níveis.
  \item Exporte apenas quando precisar de arquivo, análise externa ou partilha formal.
\end{enumerate}

\subsection{Leitura prática}
O Dashboard não serve para operações transacionais do dia a dia. É uma zona de leitura estratégica. Quando o problema for um colaborador concreto, o ponto certo costuma ser Colaboradores, Aprovações, Férias / Ausências ou Admissões.

\chapter{Gestão de colaboradores}
\section{Colaboradores}
Esta é a área central do registo de pessoas. Aqui RH encontra a ficha consolidada, permissões, estado e ações de gestão.

\subsection{O que RH pode fazer}
\begin{itemize}
  \item Pesquisar colaboradores por nome, email, equipa ou estado.
  \item Abrir a ficha completa de um colaborador.
  \item Editar dados quando o fluxo permitir.
  \item Ver e ajustar permissões.
  \item Ativar ou desativar contas.
  \item Exportar e importar, quando os botões estiverem disponíveis.
\end{itemize}

\subsection{Estrutura interna do detalhe}
Em muitos casos o detalhe do colaborador agrupa-se em três blocos de trabalho:
\begin{itemize}
  \item \textbf{Ficha}: dados e informação operacional.
  \item \textbf{Permissões}: acessos e poderes do utilizador.
  \item \textbf{Estado}: ativo, inativo ou contexto equivalente.
\end{itemize}

\subsection{Fluxo de validação de registo}
\begin{enumerate}
  \item Localize o colaborador.
  \item Abra a ficha.
  \item Confirme país, equipa, email, nome e estado.
  \item Corrija o que estiver incompleto.
  \item Grave e verifique se o sistema gerou feedback de sucesso.
\end{enumerate}

\subsection{Quando usar permissões}
Use a área de permissões quando o colaborador precisar de acesso adicional, quando um acesso tiver de ser revogado ou quando a revisão de poder de ação fizer parte do processo de RH.

\chapter{Aprovações}
\section{Aprovações}
É o centro de decisão para pedidos que exigem validação humana: alterações de ficha, férias, ausências e admissões.

\subsection{Separadores principais}
\begin{itemize}
  \item \textbf{Perfis}: pedidos de alteração de dados.
  \item \textbf{Férias}: pedidos de férias e ausências.
  \item \textbf{Admissões}: validação de novos colaboradores.
\end{itemize}

\subsection{Como aprovar ou rejeitar um pedido}
\begin{enumerate}
  \item Abra a fila de aprovações.
  \item Escolha o separador correto.
  \item Abra o pedido em detalhe.
  \item Leia o contexto completo, incluindo comentários, campos e anexos.
  \item Confirme o impacto operacional.
  \item Execute a decisão com a justificação necessária.
\end{enumerate}

\subsection{Boas decisões}
\begin{itemize}
  \item Não aprovar sem perceber o efeito da alteração.
  \item Não rejeitar sem explicar o motivo quando o fluxo pede justificação.
  \item Confirmar sempre se o pedido corresponde ao colaborador certo.
\end{itemize}

\chapter{Admissões}
\section{Admissões}
O módulo de admissões acompanha o onboarding e a validação de novos colaboradores.

\subsection{O que pode aparecer}
\begin{itemize}
  \item Lista de admissões em curso.
  \item Filtros por estado.
  \item Detalhe com dados pessoais, documentos e contrato.
  \item Ações para aprovar, corrigir ou concluir.
\end{itemize}

\subsection{Fluxo recomendado de análise}
\begin{enumerate}
  \item Abrir a admissão.
  \item Rever dados pessoais primeiro.
  \item Passar depois aos documentos.
  \item Validar contrato e campos finais.
  \item Aprovar os dados corretos ou rejeitar os que precisam de correção.
  \item Concluir a admissão quando a informação estiver fechada.
\end{enumerate}

\subsection{Interpretação dos estados}
\begin{itemize}
  \item \textbf{Aguarda análise}: o pedido está pendente.
  \item \textbf{Correção pedida}: existem campos a corrigir.
  \item \textbf{Aguarda contrato}: falta etapa contratual.
  \item \textbf{Concluído}: o processo está fechado.
\end{itemize}

\chapter{Banco de Horas}
\section{Banco de Horas}
O Banco de Horas serve para acompanhar créditos, débitos e saldos. É um módulo sensível porque qualquer movimento altera o histórico do colaborador.

\subsection{Separadores típicos}
\begin{itemize}
  \item \textbf{Meu Saldo}: leitura do colaborador.
  \item \textbf{Visão RH}: visão de gestão.
  \item \textbf{Lançamentos}: movimentos individuais.
  \item \textbf{Limites}: configuração e regras.
  \item \textbf{Relatórios}: exportações e documentos de apoio.
\end{itemize}

\subsection{Fluxo de trabalho}
\begin{enumerate}
  \item Abra a página.
  \item Escolha o separador adequado.
  \item Localize a pessoa e o período.
  \item Adicione ou consulte o movimento necessário.
  \item Valide a origem e o motivo.
  \item Guarde e confirme o efeito no saldo.
\end{enumerate}

\subsection{Cuidados essenciais}
\begin{itemize}
  \item Registar sempre o motivo correto.
  \item Evitar movimentos sem validação documental quando isso for exigido.
  \item Confirmar se o lançamento foi feito para a pessoa certa.
\end{itemize}

\chapter{Férias e ausências}
\section{Férias / Ausências}
Esta é uma área crítica para RH porque concentra o planeamento, a leitura operacional e a manutenção da política de ausências.

\subsection{Separadores principais e respetivo uso}
\begin{itemize}
  \item \textbf{Resumo}: visão geral do ano e saldos.
  \item \textbf{Calendário}: leitura diária e visualização temporal.
  \item \textbf{Férias da equipa}: planeamento coletivo e análise de conflitos.
  \item \textbf{Dias automáticos}: manutenção da política de dias dados pela empresa.
  \item \textbf{Mapa de Férias}: exportação e consulta consolidada.
  \item \textbf{Atribuição imediata}: operações rápidas, quando permitidas.
\end{itemize}

\subsection{Resumo}
O resumo é o ecrã de leitura imediata. Mostra o essencial para perceber se o colaborador tem dias disponíveis, pedidos pendentes ou consumos já registados.

\subsection{Calendário}
O calendário é a melhor vista para cruzar férias, ausências, feriados e dias automáticos. RH deve usá-lo antes de aprovar pedidos com risco de sobreposição.

\subsection{Férias da equipa}
Este separador ajuda a responder a perguntas como: quem está fora, quem já pediu férias, e em que período a equipa fica mais frágil.

\subsection{Mapa de Férias}
É a vista adequada quando é necessário exportar ou partilhar um consolidado para análise superior ou planeamento externo ao portal.

\subsection{Atribuição imediata}
Quando ativada, esta funcionalidade acelera o registo direto de um período sem passar por todo o fluxo tradicional.

\subsection{Dias automáticos}
Os dias automáticos registam dias dados pela empresa por ano e por âmbito. A regra é particularmente importante em organizações com Portugal e Brasil, porque o que é comum pode existir nos dois países e o que é local pode ser específico de um deles.

\subsubsection{O que RH pode fazer}
\begin{itemize}
  \item Definir o ano.
  \item Escolher o âmbito geográfico.
  \item Criar, remover ou corrigir dias.
  \item Confirmar se a alteração ficou refletida no calendário.
\end{itemize}

\subsubsection{Fluxo recomendado}
\begin{enumerate}
  \item Entre em \textit{Férias / Ausências}.
  \item Abra \textit{Dias automáticos}.
  \item Selecione o ano alvo.
  \item Verifique os dias já existentes.
  \item Adicione, altere ou remova conforme a política vigente.
  \item Grave e confirme o resultado na lista e no calendário.
\end{enumerate}

\subsection{Regras operacionais}
\begin{itemize}
  \item Um pedido deve ser lido no contexto da equipa.
  \item Uma ausência médica pode ter regras diferentes de férias.
  \item Dias automáticos não devem ser tratados como pedidos normais de colaborador.
\end{itemize}

\chapter{Formações}
\section{Formações}
Este módulo serve para consulta, atribuição, acompanhamento e fecho de formações.

\subsection{O que RH vê na prática}
\begin{itemize}
  \item Lista de formações com pesquisa.
  \item Filtros avançados por estado, origem, entidade, horas e datas.
  \item Botão de exportação para Excel.
  \item Ação de atribuir a um ou vários colaboradores, quando permitido.
  \item Botão de concluir quando o processo o admitir.
\end{itemize}

\subsection{Fluxo de atribuição}
\begin{enumerate}
  \item Abrir \textit{Formações}.
  \item Escolher os colaboradores.
  \item Preencher nome, link, horas, data de início e entidade.
  \item Confirmar se a formação é própria ou atribuída.
  \item Guardar e validar se o registo ficou visível na lista.
\end{enumerate}

\subsection{Fluxo de conclusão}
\begin{enumerate}
  \item Abrir a formação.
  \item Confirmar se a pessoa correta está associada.
  \item Verificar se o conteúdo foi de facto realizado.
  \item Marcar como concluída.
\end{enumerate}

\chapter{Equipas}
\section{Equipas}
A página de Equipas é útil para RH perceber estrutura, relações e distribuição interna.

\subsection{Quando usar}
\begin{itemize}
  \item Confirmar a composição de uma equipa.
  \item Ver lideranças e coordenações.
  \item Cruzar a estrutura com férias, formações e admissões.
\end{itemize}

\chapter{Saúde e bem-estar}
\section{Saúde e bem-estar}
Para RH, esta área é tanto um repositório de conteúdos como uma superfície de administração de informação sensível.

\subsection{Conteúdos comuns nas duas geografias}
\begin{itemize}
  \item Formulário de assédio.
  \item Ergonomia.
  \item Suporte Básico de Vida.
\end{itemize}

\subsection{O que RH pode fazer}
\begin{itemize}
  \item Editar texto, descrição e botão de cada bloco.
  \item Adicionar e remover blocos.
  \item Carregar PDFs.
  \item Ajustar a ordem ou o conteúdo visível por país, quando aplicável.
\end{itemize}

\subsection{Como pensar a página}
\begin{itemize}
  \item O que é transversal deve manter-se igual nos dois países.
  \item O que é específico deve ser isolado por país.
  \item O formulário de assédio deve ser tratado com máxima clareza e cuidado operacional.
\end{itemize}

\subsection{Fluxo de edição}
\begin{enumerate}
  \item Abrir \textit{Saúde e bem-estar}.
  \item Entrar em modo de edição.
  \item Alterar títulos, descrições e botões.
  \item Carregar PDFs quando necessário.
  \item Guardar e verificar se a versão publicada ficou correta.
\end{enumerate}

\chapter{Notificações}
\section{Notificações}
RH deve usar as notificações como fila rápida de leitura e resposta.

\subsection{Boa rotina}
\begin{itemize}
  \item Ler notificações logo no início do dia.
  \item Abrir os pedidos pendentes pela própria notificação.
  \item Marcar como lido o que já foi tratado.
\end{itemize}

\chapter{Fluxos detalhados de trabalho}
\section{Validação de dados de colaborador}
\begin{enumerate}
  \item Procurar o colaborador.
  \item Abrir a ficha.
  \item Verificar dados-base e estado.
  \item Corrigir o que estiver incorreto.
  \item Confirmar se existe fluxo de aprovação após a edição.
\end{enumerate}

\section{Aprovação de férias}
\begin{enumerate}
  \item Abrir a aprovação.
  \item Ler datas, tipo e contexto.
  \item Cruzar com a disponibilidade da equipa.
  \item Aprovar ou rejeitar com justificação.
\end{enumerate}

\section{Tratamento de admissão}
\begin{enumerate}
  \item Abrir a admissão pendente.
  \item Rever dados pessoais, documentos e contrato.
  \item Pedir correção do que estiver incompleto.
  \item Concluir apenas quando o processo estiver fechado.
\end{enumerate}

\section{Atualização da página de bem-estar}
\begin{enumerate}
  \item Abrir a página.
  \item Decidir se o conteúdo é comum ou específico de país.
  \item Editar os blocos.
  \item Confirmar que os blocos comuns ficaram iguais em PT e BR.
  \item Guardar e testar o resultado final.
\end{enumerate}

\chapter{Responsabilidades e controlo}
\section{Segurança de uso}
\begin{itemize}
  \item Trabalhe sempre dentro do âmbito da sua permissão.
  \item Confirme o país e o contexto antes de editar.
  \item Registe a decisão com clareza para que outra pessoa consiga seguir o histórico.
\end{itemize}

\section{Erros comuns}
\begin{itemize}
  \item Abrir o colaborador errado.
  \item Editar o país errado.
  \item Aprovar um pedido sem ler o contexto da equipa.
  \item Tratar um bloco comum como se fosse um bloco local.
\end{itemize}

\chapter{Anexo de ecrãs detalhados}
\section{Colaboradores}
Este é o ecrã que RH vai usar mais vezes para consultas e correções. A sua utilidade está na combinação entre pesquisa rápida, detalhe rico e ações de gestão.

\subsection{Antes de editar}
\begin{itemize}
  \item Confirme sempre o utilizador correto.
  \item Veja o estado da conta.
  \item Leia a equipa, o país e o papel antes de avançar.
\end{itemize}

\subsection{Depois de editar}
\begin{itemize}
  \item Confirme se o sistema indicou sucesso.
  \item Verifique se a alteração refletiu no detalhe.
  \item Considere se a ação exige notificação ou validação posterior.
\end{itemize}

\section{Detalhe do colaborador}
O detalhe do colaborador funciona como centro de comando. A partir daqui RH pode trabalhar com a ficha, com as permissões e com o estado da conta.

\subsection{Três perguntas que RH deve fazer sempre}
\begin{enumerate}
  \item Quem é a pessoa certa?
  \item Qual é o contexto da alteração?
  \item Esta alteração é apenas edição ou é um pedido que exige aprovação?
\end{enumerate}

\subsection{Leitura por abas}
\begin{itemize}
  \item \textbf{Ficha}: dados estruturais e administrativos.
  \item \textbf{Permissões}: acessos e poderes de execução.
  \item \textbf{Estado}: ativo/inativo e ciclo da conta.
\end{itemize}

\section{Dashboard}
\subsection{Interpretação correta}
O Dashboard não substitui a ficha do colaborador. Ele ajuda a responder a perguntas mais amplas: onde estão os maiores grupos, que distribuição existe, que tendências aparecem e onde pode haver risco ou oportunidade.

\subsection{Uso avançado}
\begin{itemize}
  \item Drill-down para ir do geral para o detalhe.
  \item Roll-up para voltar ao nível anterior.
  \item Drill-through para exportação ou detalhe fora do gráfico.
\end{itemize}

\section{Aprovações}
\subsection{Fluxos de decisão}
\begin{itemize}
  \item Aprovação total quando a informação está correta.
  \item Rejeição total quando o pedido não pode avançar.
  \item Rejeição parcial quando apenas uma parte do registo está incorreta.
\end{itemize}

\subsection{Boa disciplina}
\begin{enumerate}
  \item Ler.
  \item Confirmar.
  \item Decidir.
  \item Registar a razão.
\end{enumerate}

\section{Admissões}
\subsection{Leitura do ecrã}
As admissões pedem disciplina porque misturam documentos, dados pessoais e contrato. Em contexto RH, a pressa é normalmente o maior risco.

\subsection{O que validar primeiro}
\begin{itemize}
  \item Nome e dados-base.
  \item Documentos anexados.
  \item Informação contratual.
\end{itemize}

\section{Banco de Horas}
\subsection{Como evitar erros}
\begin{itemize}
  \item Nunca lançar sem motivo claro.
  \item Nunca alterar o saldo sem perceber a origem.
  \item Nunca assumir que o movimento é para a pessoa certa sem confirmação do detalhe.
\end{itemize}

\section{Férias / Ausências}
\subsection{Leitura analítica}
RH deve usar o módulo com três perspetivas simultâneas: individual, equipa e calendário global. Só assim se percebe se um pedido é apenas um pedido ou um risco operacional.

\subsection{Dias automáticos}
\begin{itemize}
  \item São uma regra da empresa.
  \item Podem ser comuns às duas geografias.
  \item Devem ser mantidos coerentes com o calendário e com a política interna.
\end{itemize}

\section{Formações}
\subsection{Atribuição responsável}
Quando RH atribui formação, precisa garantir três coisas: a formação é relevante, a pessoa certa foi escolhida e o registo ficará fácil de acompanhar depois.

\subsection{Conclusão de formação}
\begin{itemize}
  \item Verificar se a formação foi realizada.
  \item Confirmar se a conclusão precisa de evidência adicional.
  \item Marcar como concluída apenas quando o processo estiver completo.
\end{itemize}

\section{Equipas}
\subsection{Quando RH usa Equipas}
\begin{itemize}
  \item Para confirmar hierarquia.
  \item Para validar líderes e membros.
  \item Para ligar estrutura a férias, ausências, formações e admissões.
\end{itemize}

\section{Saúde e bem-estar}
\subsection{Governança de conteúdo}
\begin{itemize}
  \item Conteúdo comum deve permanecer comum.
  \item Conteúdo local deve ficar isolado por país.
  \item A página deve ser legível sem criar dúvidas sobre o que é transversal e o que é específico.
\end{itemize}

\subsection{Fluxo de manutenção recomendado}
\begin{enumerate}
  \item Ler o conteúdo atual.
  \item Identificar o que é comum e o que é local.
  \item Editar com a equipa certa.
  \item Publicar e testar no país certo.
\end{enumerate}

\section{Notificações}
\subsection{Como RH deve usá-las}
\begin{itemize}
  \item Como fila de trabalho do dia.
  \item Como entrada para pedidos pendentes.
  \item Como memória curta do que foi tratado ou ficou em espera.
\end{itemize}

\chapter{Guia de nomenclatura prática}
\section{PT-PT e PT-BR no trabalho diário}
\begin{longtable}{|p{0.3\textwidth}|p{0.6\textwidth}|}
\hline
\textbf{PT-PT} & \textbf{PT-BR / observação} \\
\hline
Colaboradores & Colaboradores / pessoas / usuários internos, dependendo do contexto. \\
\hline
Chefes de equipa & Líderes de equipe / gestores de equipe. \\
\hline
A Minha Ficha & Meu cadastro / meu perfil. \\
\hline
Férias / Ausências & Mantém-se, mas a leitura administrativa pode variar por país. \\
\hline
Formações & Treinamentos. \\
\hline
Banco de Horas & Mantém-se o conceito, embora a linguagem de apoio possa variar. \\
\hline
\end{longtable}

\chapter{Conclusão}
Este manual deve ser atualizado sempre que houver mudanças de processo, novos separadores, novos fluxos de aprovação ou mudanças de nomenclatura entre PT-PT e PT-BR. Em RH, a documentação só é útil se acompanhar a plataforma de forma contínua.

\chapter{Anexo operacional por ecrã}
\section{Mapa rápido dos ecrãs de RH}
\begin{longtable}{|p{0.18\textwidth}|p{0.14\textwidth}|p{0.17\textwidth}|p{0.21\textwidth}|p{0.24\textwidth}|}
\hline
\textbf{Ecrã} & \textbf{Quem usa} & \textbf{Abertura} & \textbf{Botões / abas-chave} & \textbf{Notas PT-PT / PT-BR} \\
\hline
Dashboard & Root / access total & \texttt{/dashboard} & Drill, filtros, exportação & É um ecrã executivo; no Brasil a leitura de geografia e função pode ser mais detalhada. \\
\hline
Colaboradores & RH / gestão & \texttt{/colaboradores} & Ficha, permissões, estado, importar, exportar & PT-PT usa muito ficha; PT-BR tende a usar cadastro/perfil na conversa. \\
\hline
Aprovações & RH / chefias autorizadas & \texttt{/aprovacoes} & Profiles, vacations, admissions, aprovar, rejeitar & O significado funcional é o mesmo, mas a legenda pode variar entre países. \\
\hline
Admissões & RH & \texttt{/admissoes} & Filtros, detalhe, aprovar dados pessoais, pedir correção & Campos documentais diferem bastante entre PT e BR. \\
\hline
Banco de Horas & RH / autorizados & \texttt{/banco-horas} & Meu Saldo, Visão RH, Lançamentos, Limites, Relatórios & Em PT-BR o uso de saldo e relatórios tende a ser a linguagem mais natural. \\
\hline
Férias / Ausências & RH / chefias & \texttt{/ferias} & Resumo, Calendário, Férias da equipa, Dias automáticos, Mapa, Atribuição imediata & BR pode expor âmbitos como BR\_SP ou BR\_RS; PT usa focos de política local. \\
\hline
Formações & RH / autorizados & \texttt{/formacoes} & Nova formação, Exportar Excel, concluir, filtros & Treinamentos em PT-BR é equivalente funcional; a UI mantém Formações. \\
\hline
Equipas & RH / chefias & \texttt{/equipas} & Modal de detalhe, comparação, visão de membros & A leitura de equipe costuma ser mais intuitiva em PT-BR; a UI conserva o termo equipa. \\
\hline
Saúde e bem-estar & RH / editor autorizado & \texttt{/saude-bem-estar} & Editar página, adicionar blocos, carregar PDFs, Reportar situação & Formulário de assédio, ergonomia e suporte básico de vida são comuns às duas geografias. \\
\hline
Notificações & Todos & \texttt{/notifications} & Ler, filtrar, marcar, apagar, abrir links & Mensagens técnicas devem ser humanizadas ao orientar o utilizador. \\
\hline
\end{longtable}

\section{Dashboard em detalhe}
\subsection{Objetivo do ecrã}
O Dashboard responde a perguntas macro sobre a organização. Não deve ser lido como ficha de pessoa nem como área de decisão operacional do dia a dia.

\subsection{Botões e ações}
\begin{itemize}
  \item Drill-up: voltar a um nível superior.
  \item Roll-up: compactar ou regressar à visão agregada.
  \item Drill-through: detalhar e exportar o ponto selecionado.
  \item Exportação: gerar Excel quando o contexto exigir análise externa.
\end{itemize}

\subsection{Exemplo de uso}
O RH quer perceber a distribuição por função e geografia. Aplica filtros, abre um gráfico, faz drill-down e exporta apenas o subconjunto necessário.

\subsection{Erro típico}
Usar o Dashboard para confirmar um dado individual que deveria ser revisto na ficha do colaborador.

\section{Colaboradores em detalhe}
\subsection{O que este ecrã resolve}
É a base de manutenção da plataforma: dados pessoais, estado, equipas, funções e acessos.

\subsection{Campos e áreas mais relevantes}
\begin{itemize}
  \item Ficha do colaborador.
  \item Permissões.
  \item Estado da conta.
  \item Associação à equipa.
  \item Exportação/importação quando aplicável.
\end{itemize}

\subsection{Diferenças por país}
\begin{itemize}
  \item PT usa dados como NIF, NISS, cartão de cidadão e IRS.
  \item BR usa CPF, PIS, CTPS, RG, CNH e título de eleitor.
  \item O ecrã é o mesmo, mas o conjunto de campos visíveis pode variar pelo país e pelo contexto da pessoa.
\end{itemize}

\subsection{Exemplo de uso}
Uma colaboradora muda de equipa e de cargo. RH localiza a pessoa, confirma o país, altera a ficha e ajusta permissões se necessário.

\subsection{Erro típico}
Editar uma ficha sem confirmar se o utilizador está em Portugal ou no Brasil, o que pode levar a alterações no campo errado.

\section{Aprovações em detalhe}
\subsection{Três filas de trabalho}
\begin{itemize}
  \item Perfis: alterações da ficha.
  \item Férias: validação de ausências.
  \item Admissões: entrada de novos colaboradores.
\end{itemize}

\subsection{Fluxo correto}
\begin{enumerate}
  \item Ler o item pendente.
  \item Confirmar a natureza do pedido.
  \item Verificar o contexto.
  \item Tomar decisão.
  \item Registar a justificação quando for obrigatório.
\end{enumerate}

\subsection{Exemplo de uso}
Uma alteração de ficha chega à fila. RH vê a notificação, abre a aprovação, compara os campos antigos e novos e decide se aprova, rejeita ou pede correção.

\subsection{Erro típico}
Rejeitar sem explicar o motivo. Isto gera retrabalho e atraso no processo.

\section{Admissões em detalhe}
\subsection{Leitura da página}
\begin{itemize}
  \item Lista de admissões com pesquisa e filtro.
  \item Estatuto da admissão.
  \item Detalhe com dados pessoais, documentos e contrato.
\end{itemize}

\subsection{Diferenças PT-PT / PT-BR}
\begin{itemize}
  \item Documentos aceites ou esperados mudam bastante.
  \item Alguns campos fiscais são específicos de cada país.
  \item O estágio de contrato pode ter nomenclaturas internas ligeiramente diferentes, mas a função da página é a mesma.
\end{itemize}

\subsection{Exemplo de uso}
RH abre uma admissão em estado 	extit{Aguarda contrato}, valida os campos já recebidos e pede a conclusão da documentação contratual.

\subsection{Erro típico}
Concluir uma admissão sem confirmar documentos obrigatórios.

\section{Banco de Horas em detalhe}
\subsection{Leitura dos separadores}
\begin{itemize}
  \item Meu Saldo: leitura individual.
  \item Visão RH: leitura agregada.
  \item Lançamentos: movimentos.
  \item Limites: regras.
  \item Relatórios: documentação periódica.
\end{itemize}

\subsection{Diferenças por perfil e país}
\begin{itemize}
  \item Em PT-BR a leitura de relatórios pode exigir mais contexto de geografia.
  \item Em PT-PT pode haver maior peso na interpretação do saldo e da política interna.
  \item O ecrã adapta-se ao utilizador, mas o conceito base é o mesmo.
\end{itemize}

\subsection{Exemplo de uso}
RH precisa de registar um crédito de horas. Seleciona a pessoa, escolhe o tipo, escreve o motivo, confirma o valor e grava o lançamento.

\subsection{Erro típico}
Lançar horas no colaborador errado ou sem descrever o motivo de forma verificável.

\section{Férias / Ausências em detalhe}
\subsection{Ler o módulo como um sistema, não como uma lista}
O módulo junta política, calendário, equipa e aprovação. RH deve lê-lo de forma integrada.

\subsection{Dias automáticos e diferenças entre países}
\begin{itemize}
  \item PT: normalmente mais centrado em férias, tolerâncias e dias da empresa localizados no contexto português.
  \item BR: pode incluir âmbitos por estado, como SP ou RS, e regras específicas por geografia.
  \item O bloco comum deve permanecer idêntico nas duas geografias quando a política for transversal.
\end{itemize}

\subsection{Exemplo de uso}
RH cria um dia automático de empresa para o ano em curso, valida o âmbito e confirma se aparece no calendário certo.

\subsection{Erro típico}
Introduzir um dia automático como se fosse uma ausência individual.

\section{Formações em detalhe}
\subsection{O que o ecrã permite}
\begin{itemize}
  \item Atribuir formações.
  \item Acompanhar estado.
  \item Procurar por colaborador, origem ou entidade.
  \item Exportar a lista.
\end{itemize}

\subsection{Diferenças de linguagem}
\begin{itemize}
  \item PT-PT: Formações, conclusão, entidade.
  \item PT-BR: o mesmo conceito pode ser lido como treinamentos, mas a interface mantém a estrutura de RH.
\end{itemize}

\subsection{Exemplo de uso}
RH atribui uma formação, revê o estado mais tarde e conclui o registo quando a pessoa terminar o conteúdo.

\subsection{Erro típico}
Criar formações duplicadas porque não foi usada a pesquisa antes da criação.

\section{Equipas em detalhe}
\subsection{Objetivo operacional}
Perceber estrutura, liderança e distribuição dos membros para apoiar férias, admissões e formações.

\subsection{Exemplo de uso}
RH abre uma equipa, revê os membros e valida se a chefia atribuída corresponde ao organograma esperado.

\subsection{Erro típico}
Usar a página como um diretório sem cruzar com outros módulos.

\section{Saúde e bem-estar em detalhe}
\subsection{Governança por país}
\begin{itemize}
  \item Conteúdos comuns: formulário de assédio, ergonomia, suporte básico de vida.
  \item Conteúdos locais: documentos e políticas próprios de PT ou BR.
\end{itemize}

\subsection{Exemplo de uso}
RH publica o mesmo formulário de assédio nas duas geografias, mas adapta o texto de apoio local quando necessário.

\subsection{Erro típico}
Misturar conteúdo transversal com conteúdo específico e gerar duplicação desnecessária.

\section{Notificações em detalhe}
\subsection{Uso operacional}
As notificações são a primeira fila de triagem do RH. O ideal é tratá-las cedo para evitar acumulação.

\subsection{Exemplo de uso}
Uma notificação de admissão chega com link para aprovação. RH abre, valida e segue o fluxo sem sair da mensagem.

\subsection{Erro típico}
Ler a notificação e esquecer de clicar na ação associada.

\chapter{Cenários de uso avançado}
\section{Cenário 1: alteração de ficha com aprovação}
\begin{enumerate}
  \item O colaborador submete uma alteração.
  \item RH recebe notificação.
  \item RH abre Aprovações > Perfis.
  \item Compara antes/depois.
  \item Aprova ou rejeita.
\end{enumerate}

\section{Cenário 2: férias com risco de sobreposição}
\begin{enumerate}
  \item O colaborador pede férias.
  \item RH cruza com o calendário da equipa.
  \item Verifica dias automáticos e feriados.
  \item Decide com base no impacto operacional.
\end{enumerate}

\section{Cenário 3: conteúdo de bem-estar transversal}
\begin{enumerate}
  \item RH abre Saúde e bem-estar.
  \item Confirma que o bloco é comum.
  \item Publica o mesmo conteúdo em PT e BR.
  \item Testa a visualização em ambas as geografias.
\end{enumerate}

\chapter{Checklist de qualidade}
\begin{itemize}
  \item O ecrã certo foi aberto?
  \item O país certo foi considerado?
  \item O conteúdo é comum ou local?
  \item A ação foi gravada e confirmada?
  \item A notificação foi lida e o fluxo seguiu em frente?
\end{itemize}

\chapter{RH Brasil: operação local aprofundada}
\section{Foco cadastral e documental no Brasil}
No contexto de RH Brasil, estes pontos são críticos em validação:
\begin{itemize}
  \item CPF, PIS e CTPS (incluindo série e data de expedição).
  \item RG (órgão emissor e data de expedição).
  \item CNH (categoria e validade), quando aplicável.
  \item Título de eleitor, zona e seção.
  \item Estado de trabalho (SP/RS/outro escopo BR configurado).
\end{itemize}

\section{Matriz de decisão em aprovações (BR)}
\begin{longtable}{|p{0.24\textwidth}|p{0.24\textwidth}|p{0.2\textwidth}|p{0.24\textwidth}|}
\hline
\textbf{Tipo de solicitação} & \textbf{Condição mínima} & \textbf{Ação} & \textbf{Motivo recomendado} \\
\hline
Alteração cadastral & CPF/RG/CTPS consistentes & Aprovar & Cadastro validado \\
\hline
Alteração cadastral & Documento incompleto ou inconsistente & Rejeitar / solicitar correção & Informar campo exato \\
\hline
Férias & Cobertura operacional preservada & Aprovar & Planejamento viável \\
\hline
Férias & Sobreposição crítica de equipe & Rejeitar / negociar data & Risco operacional \\
\hline
Admissão & Dados e contrato finalizados & Concluir & Onboarding completo \\
\hline
\end{longtable}

\section{Fluxos exclusivos com exemplos (BR)}
\subsection{Fluxo BR-1: validação de documentos de admissão}
\begin{enumerate}
  \item Abrir Admissões.
  \item Selecionar o candidato pendente.
  \item Revisar CPF, RG, CTPS e dados contratuais.
  \item Solicitar ajuste quando faltar evidência.
  \item Concluir quando tudo estiver fechado.
\end{enumerate}

\subsection{Fluxo BR-2: gestão de dias automáticos por escopo}
\begin{enumerate}
  \item Abrir Férias / Ausências.
  \item Ir ao separador Dias automáticos.
  \item Escolher ano e escopo (Brasil geral ou estado, quando configurado).
  \item Ajustar os dias.
  \item Salvar e validar no calendário.
\end{enumerate}

\section{Erros críticos a evitar no BR}
\begin{itemize}
  \item Registrar alterações sem conferir documentos-base.
  \item Tratar escopo estadual como se fosse escopo nacional.
  \item Fechar admissões com pendências documentais.
\end{itemize}
